\documentclass[a4paper,11pt]{article}
\usepackage[utf8]{inputenc}
\usepackage[T1]{fontenc}
\usepackage[french]{babel}
\usepackage[left=2cm,right=2cm,top=2cm,bottom=2cm]{geometry}
\usepackage{xcolor}
\usepackage{hyperref}
\usepackage{fontawesome5}
\usepackage{parskip}

% --- Couleurs Nokia ---
\definecolor{primary}{RGB}{18, 65, 145}

\hypersetup{colorlinks=true, linkcolor=primary, urlcolor=primary}

\begin{document}

% --- EN-TÊTE ---
\begin{center}
    {\Large \textbf{Loïc Aron MBASSI EWOLO}}\\[4pt]
    \small
    \faMapMarker\ Cergy, France \quad
    \faPhone\ (+33) 7 59 52 33 98 \quad
    \faEnvelope\ \href{mailto:loicaron.mbassiewolo@ensea.fr}{loicaron.mbassiewolo@ensea.fr}
\end{center}
\vspace{0.5cm}

% --- DESTINATAIRE ---
\hfill
\begin{minipage}[t]{0.45\textwidth}
    \textbf{NOKIA}\\
    Service Recrutement\\
    Lannion, France
\end{minipage}

\vspace{0.8cm}

Cergy, le 23 décembre 2025

\vspace{0.5cm}

\textbf{Objet :} Candidature Stage — Développement Logiciels Temps Réel 5G/6G\\

Madame, Monsieur,

Votre offre de stage au sein de l'équipe Mobile Networks de Nokia Lannion a immédiatement retenu mon attention. Le développement de logiciels temps réel pour les stations de base 5G/6G représente exactement le type de défi technique que je recherche : allier rigueur système, contraintes temps réel et intégration continue dans un environnement Agile.

Élève-ingénieur en double diplôme à l'ENSEA, spécialisé en Traitement du Signal et Intelligence Artificielle, j'ai acquis une solide expérience en programmation système. Mon projet PROJET-TAXI, réalisé en C sous Linux, m'a confronté aux problématiques que vous mentionnez : j'ai implémenté une architecture multiprocessus avec communication inter-processus (IPC, signaux UNIX, mémoire partagée), un ordonnanceur FIFO avec prévention des deadlocks, et une visualisation temps réel en OpenGL. Ce projet m'a appris à débugger du code concurrent et à profiler des applications bas-niveau.

J'ai également participé au Challenge CoHoMa (Défense) où j'ai intégré et optimisé des modèles d'IA sur Nvidia Jetson pour des robots autonomes. Ce travail m'a familiarisé avec le déploiement Docker sur systèmes embarqués, le traitement de données Lidar 3D et les contraintes de ressources limitées. Mon stage au laboratoire IoT de l'ENSPY a renforcé cette expertise : j'y ai optimisé des modèles TensorFlow Lite pour du hardware contraint avec des exigences temps réel strictes.

Côté méthodologie, mon expérience chez CSC Fintech m'a formé à l'intégration continue (CI/CD), aux tests unitaires et d'intégration (TDD), ainsi qu'aux revues de code en équipe Agile/Scrum. Cette rigueur dans le processus de développement correspond à l'approche que vous décrivez pour vos environnements de test et de debug.

Je suis disponible dès \textbf{avril 2026} pour une durée de 4 à 6 mois. L'opportunité de contribuer aux technologies 5G/6G chez Nokia, dans une équipe multiculturelle à Lannion, m'enthousiasme particulièrement. Je serais ravi de vous présenter mes projets lors d'un entretien.

Cordialement,

\vspace{0.8cm}

\textbf{Loïc Aron MBASSI EWOLO}

\end{document}
